\chapter{Compile-Time Generation}
\label{ch:compile-time-generation}
\chapterquote{Generated code should be boring enough to inspect and useful enough to remove repetition.}{The generation rule}

\begin{chaptermap}
Core question & What does Corusco generate, why is it generated at compile time, and how should application code consume the generated companions without losing design control?\\
Primary annotations & \api{@SwingForm}, \api{@TextField}, \api{@DateField}, \api{@ComboBox}, \api{@CheckBox}, \api{@SwingTable}, \api{@Column}, \api{@UiAction}, validation annotations, and \api{@Help}.\\
Generated outputs & Field keys, resource keys, problem codes, descriptors, form models, view contracts, behavior plans, binding facades, table descriptors, command descriptors, command factories, and deterministic companion names.\\
Build topic & Annotation-processor configuration, generated-source inspection, processor diagnostics, generated-source compiler tests, modular builds, and stale-output hygiene.\\
Companion example & \generationExample.\\
\end{chaptermap}

\section{Generation is a compile-time agreement}

Corusco generation starts from an ordinary Java source declaration and ends in
ordinary Java source. A record annotated with \api{@SwingForm} describes an
editable value shape. A record annotated with \api{@SwingTable} describes a
row shape and its table columns. A presenter method annotated with
\api{@UiAction} describes a user operation. The processor reads those
declarations during javac annotation processing, validates them, and writes
generated companions beside the annotated source type. Runtime code imports
the generated classes. It does not scan annotations reflectively after the
application has already started.

\termdef{Compile-time generation}{in Corusco} means that javac is part of the
UI metadata pipeline. The compiler sees the annotated source. The processor
sees the compiler model of that source through \api{javax.lang.model}. The
processor writes Java files. The compiler then compiles those Java files with
the rest of the program. The generated classes can be opened in the IDE,
searched with normal code search, checked by tests, packaged by Gradle, and
included in a runtime image. This is a very different contract from a runtime
metadata scanner that discovers annotations after startup.

The agreement has three sides. The annotated source promises stable facts:
field ids, column ids, action ids, resource ids, validation constraints, help
topics, and supported component shapes. The processor promises deterministic
generated Java and useful diagnostics when the promise is broken. Application
code promises to consume the generated companions as ordinary code, not to
reconstruct ids by string convention and not to edit generated files by hand.
When all three sides are kept, generation removes mechanical work without
making the screen mysterious.

\termdef{Generated companion}{in this chapter} is a generated class or
interface whose name is derived from the annotated source. A form record named
\api{CustomerEdit} may produce \api{CustomerEditFields},
\api{CustomerEditResources}, \api{CustomerEditProblems},
\api{CustomerEditDescriptors}, \api{CustomerEditFormModel},
\api{CustomerEditView}, \api{CustomerEditBehaviorPlan}, and
\api{CustomerEditBindings}. A table row record named \api{CustomerRow} may
produce \api{CustomerRowColumns}, \api{CustomerRowTableResources},
\api{CustomerRowTableDescriptor}, and \api{CustomerRowTableBindings}. A
presenter class named \api{CustomerPresenter} may produce
\api{CustomerPresenterActions}. The names are long enough to say what they are
and regular enough to be guessed during review.

\begin{principlebox}{Generated Means Inspectable}
Generated code may be repetitive. It should not be mysterious. If a binding,
descriptor, key, or command factory is surprising, open the generated source
and read it like Java.
\end{principlebox}

\begin{figure}[htbp]
\centering
\begin{minipage}[t]{0.48\linewidth}
\includegraphics[width=\linewidth]{\behaviorPlanFigure}
\end{minipage}\hfill
\begin{minipage}[t]{0.48\linewidth}
\includegraphics[width=\linewidth]{\coruscoTableDescriptorFigure}
\end{minipage}
\caption{Generated companions become visible through behavior plans and descriptor-backed tables. The generated Java bridges annotated source to ordinary Swing screens.}
\label{fig:generation-visual-companions}
\end{figure}

Figure~\ref{fig:generation-visual-companions} shows two places where generated
facts become visible. A behavior plan installs field bindings, validation
decoration, help, and focus behavior into a view. A table descriptor gives a
\api{JTable} stable columns, resource keys, renderable values, and state
metadata. Neither figure implies that generation owns the whole screen. The
panel is still a panel. The table is still a table. MigLayout, FlatLaf, table
renderers, presenters, and application services remain ordinary code. The
generated part is the repeated contract that would otherwise be handwritten in
several places.

This chapter is therefore not a chapter about a GUI builder. A GUI builder
tries to own component placement and often stores design in a special format.
Corusco generation owns metadata and routine wiring. Layout remains visible in
Swing code. Workflow remains visible in presenters and services. Business
rules remain visible in domain code. If a generated class contains a surprising
business decision, either the generator is doing too much or the application is
placing the wrong fact in an annotation.

Generation is also not dependency injection. It does not discover services,
build repositories, or decide which task executor a presenter should use. It
can generate command descriptors for presenter methods, but the presenter is
still the object that knows how to save, delete, refresh, or export. It can
generate form validation rules for supported field constraints, but
cross-record business validation remains Java code. The boundary should feel
conservative. Generation is useful because it removes repetition from boring
places, not because it promises to design the application.

For a reader new to Corusco, generated source is often the shortest tutorial.
Open an annotated record. Open the generated fields companion and see typed
keys. Open the generated resources companion and see typed resource ids. Open
the form model and see \api{TextFieldModel} and \api{FieldModel} instances.
Open the behavior plan and see direct calls to Swing behavior factories. The
processor implementation is important, but generated output explains the
runtime surface that application code will actually use.

\section{Stable source declarations}

Generation is appropriate only when the source declaration already contains a
stable screen fact. A record component annotated as a text field contains a
field name, owner type, value type, and component intent. A table component
annotated as a column contains a row type, value type, column order, width
policy, resource ids, and persistence id. A no-argument presenter method
annotated as a UI action contains an operation identity, text resource id,
tooltip resource id, optional shortcut metadata, optional mnemonic, and
toggle-style flag. The processor is allowed to copy those facts into several
companions because the facts already belong to the declaration.

\termdef{Stable id grammar}{for current generated metadata} is deliberately
small. The accepted ids use letters, digits, dots, underscores, dashes, and
slashes. This is enough for resource namespaces, feature prefixes, and
human-readable diagnostics. It is not an invitation to encode the entire
screen structure into a string. Prefer \api{customer/edit/name} to
\api{leftPanelRowThreeNameTextField}. Prefer \api{customer/save} to
\api{saveButtonAction}. The id should name the UI contract, not the current
component placement.

\begin{tabularx}{0.96\linewidth}{@{}p{0.23\linewidth}p{0.24\linewidth}X@{}}
\toprule
Source kind & Stable id role & Generated uses\\
\midrule
\api{@SwingForm} & Form prefix & Field ids, resource ids, problem codes, component contracts, descriptors\\
\api{@SwingTable} & Table prefix & Table id, default column ids, table-state boundary, descriptor ownership\\
\api{@Column} & Column id & Column key, header key, tooltip key, persistence id, order, width metadata\\
\api{@UiAction} & Action id & Action key, command descriptor, resource keys, command factory, command set\\
\api{@Help} & Help and tooltip ids & Help topic metadata and optional static tooltip resource ids\\
\bottomrule
\end{tabularx}

Changing an id is a compatibility decision. Generated field keys may appear in
tests. Generated resource keys may appear in localization bundles. Generated
column persistence ids may appear in saved table layouts. Generated action ids
may appear in diagnostics, command palettes, help, or user shortcuts. A
renamed Java record component may be a refactoring; a changed generated id may
be a migration. The distinction is crucial when a screen has shipped.

Source retention is part of the design. The annotations are retained for the
compiler, not for runtime reflection. This keeps runtime modules smaller, keeps
startup simple, and gives javac the chance to reject invalid metadata while the
source element is still present. A modular application does not need to open
packages for reflective scanners merely so a UI framework can discover
annotations at runtime. The generated companions are the runtime API.

\begin{pitfallbox}{Runtime Annotation Nostalgia}
Do not write runtime code that searches for \api{@SwingForm},
\api{@SwingTable}, or \api{@UiAction}. Runtime code should use generated
companions. The annotations are compile-time inputs.
\end{pitfallbox}

The current processor supports specific source shapes. A form source is a
non-generic record annotated with \api{@SwingForm} and at least one field-kind
annotation on a record component. A table source is a non-generic record
annotated with \api{@SwingTable} and at least one \api{@Column}. A UI action
source is a no-argument \api{void} method enclosed by a type. These limits are
not arbitrary. Records give the processor stable component names, value types,
constructor shape, and component order. No-argument action methods give command
factories a direct invocation target without inventing argument injection.

Field-kind annotations are mutually exclusive. A component is a text field, a
date field, a combo box, or a check box for generated form purposes. A
component with both \api{@TextField} and \api{@CheckBox} is not a richer field;
it is an ambiguous declaration. The processor rejects it. Similar strictness
appears in table metadata: duplicate column ids are rejected, invalid widths
are rejected, and conflicting tooltip declarations are rejected. Generation
works best when invalid ambiguity is a compile error rather than a surprising
runtime choice.

Generated code should not be used to hide unstable decisions. A field label is
presentation and belongs in resources. A help topic is stable routing metadata
and may be generated from \api{@Help}. A validation constraint such as
\api{@Length(max = 80)} is stable enough for generation when it is part of the
screen contract. A database uniqueness check, a permission rule, or a
workflow-dependent warning is usually application code. If an annotation value
would need access to services, task state, current user, or network results, it
probably is not a good generation input.

The package layout reinforces this boundary. Form annotations live under
\api{cz.auderis.corusco.annotations.form}. Table annotations live under
\api{cz.auderis.corusco.annotations.table}. Validation annotations live under
\api{cz.auderis.corusco.annotations.validation}. Command annotations live
under \api{cz.auderis.corusco.annotations.command}. Help annotations live under
\api{cz.auderis.corusco.annotations.help}. Import the package that matches the
source element. Older root-package imports are not the contract; using them
should fail like any other stale source.

A source declaration should be readable without knowing the processor
implementation. The annotation names should say what kind of UI contract is
being declared. The id should say where the fact belongs. The Java type should
say what value shape is being edited, rendered, or invoked. If a reviewer must
open the processor to understand whether a field is a text editor, combo box,
table column, or command, the source declaration is not carrying enough
meaning. Generated code can be verbose; annotated source should remain crisp.

It follows that annotation attributes should not be used as a private language
for application behavior. A field annotation that says \api{@TextField} is
clear. A validation annotation that says \api{@Length(max = 80)} is clear. An
imagined annotation that says \api{@OnSave("when valid and admin and not
archived")} would be the wrong direction. The moment an attribute starts to
need boolean algebra, service lookup, permissions, or sequencing, Java code is
the better notation. Corusco's generator should make the simple declarations
worth writing, not turn complex workflow into strings.

Generated ids should be chosen as if logs will print them. In production,
diagnostics may say that \api{ActionKey[customer/save]} is disabled, that
\api{FieldKey[CustomerEdit:book/customer/edit/name]} has a problem, or that a
resource id is missing. A readable id shortens debugging. A misleading id
creates a second bug in every diagnostic. This is another reason to avoid
visual placement names. When a component moves, a log that still says
\api{left-column-field-two} becomes worse than useless. It suggests the wrong
screen location.

The default ids produced from record component names are conveniences, not
carelessness. They are good when the component name already expresses the UI
contract. A component named \api{creditLimit} under a form id
\api{customer/edit} naturally becomes \api{customer/edit/credit-limit}. If the
Java component name is a temporary abbreviation, change the annotation id
explicitly or rename the component before the screen ships. Generated defaults
are easiest to live with when Java names are chosen with the same care as API
names.

One declaration may generate several resource ids. A form field may need a
label, tooltip, accessible description, and validation messages. A table column
may need a header, tooltip, export label, and help topic. A command may need
visible text, tooltip, menu text, toolbar text, and accessible description.
The current generator emits the supported core set, and application code may
add adjacent keys for richer presentation. The important convention is that
the family remains systematic. A resource author should be able to see that
\api{customer/edit/name/label} and \api{customer/edit/name/help} belong to the
same field.

\section{Generated forms from records}

\api{@SwingForm} marks a record as a generated form source. Record components
describe editable input. The record itself is the committed value shape: the
object that can be passed to domain code after editing succeeds. The generated
form model is the mutable editing shape: it owns parse state, field models,
validation state, dirty state, and guarded result creation. This separation
matches the form chapter's central rule. Widgets are not the only place where
state lives.

\begin{lstlisting}[caption={Annotated form source used by the companion example.}]
@SwingForm(id = "book/customer/edit")
public record CustomerEdit(
        @TextField
        @Required
        @Length(max = 80)
        String name,

        @TextField
        String status) {
}
\end{lstlisting}

The generated \api{CustomerEditFields} companion contains
\api{TextFieldKey<CustomerEdit,String>} constants for \api{name} and
\api{status}. The generated resource companion contains label and tooltip
resource keys. The generated problem companion contains problem codes such as
\api{book/customer/edit/name/required} and
\api{book/customer/edit/name/length}. The generated descriptors companion
contains \api{FieldDescriptor} values in record component order. The generated
form model contains \api{TextFieldModel} instances, registered fields, a
generated rule set, and a \api{createResult} method that constructs the record
from parsed values.

\begin{lstlisting}[caption={A generated field key has stable id, owner type, and value type.}]
public static final TextFieldKey<CustomerEdit, String> NAME =
        TextFieldKey.of("book/customer/edit/name",
                CustomerEdit.class, String.class);
\end{lstlisting}

This constant is more than a convenient string. It is the field identity used
by validation, problem targets, bindings, tests, diagnostics, and resource
review. If a test enters text by the generated field key, the test does not
depend on the English label or the component's row in the panel. If a problem
points at the generated field key, focus resolution can find the correct
component through the generated view contract or a component-key mapping. The
generated key gives several subsystems one stable way to talk about the same
field.

Generated form models are plain Java objects and should be tested as such. A
test can construct \api{CustomerEditFormModel} from an original record, write
raw text into \api{model.name}, inspect parse state, inspect
\api{model.problems()}, and call \api{toResult} when valid. No Swing component
is required for that test. This is one of the strongest reasons to generate a
form model rather than bind directly to widgets. The edit contract can be
tested before the view is assembled.

The generated model also fixes a subtle ownership problem. Without a form
model, raw editor text, parsed values, parse failures, validation problems,
dirty state, and result construction tend to live in different listeners.
Generation gives them a single predictable home. Text fields own raw text and
conversion state. Field models own semantic values. The form model aggregates
problems and produces the committed record only when the current edit state is
valid. A view binding mirrors editor changes into the model; it does not
become the model.

Converters are part of that ownership. A generated \api{@TextField}
\api{BigDecimal} component should not leave every panel to decide how blank
text, parse failure, scale, and formatting are represented. The generated
model can use the standard converter path. A specialized application may still
replace or extend conversion policy, but the default screen does not need a
new anonymous document listener for every numeric field. This is the kind of
mechanical consistency generation is meant to provide.

Dirty state should be read from the model, not inferred from component events.
A generated form model starts from an original record and can compare current
field values with that original shape. The Save command can observe the form's
dirty and problem state. A Cancel workflow can ask whether there are unsaved
changes. A screenshot test can place the form in clean, dirty, invalid, and
valid states without clicking through a complete UI. If dirty state is hidden
inside one button listener, none of those workflows can share it reliably.

Touched state is a related but different concern. A field may be dirty because
data was loaded programmatically. A field may be touched because the user
visited it. Validation presentation often depends on touched state: showing
every possible error before the user edits anything can make a form feel
hostile, while hiding parse errors after user input is dishonest. Generated
form metadata gives the application stable field identities and model objects;
the presentation policy can decide how touched, dirty, and problem state are
shown.

\begin{lstlisting}[caption={Using a generated form model without Swing.}]
CustomerEdit original = new CustomerEdit("Ada", "Active");
CustomerEditFormModel model = new CustomerEditFormModel(original);

model.name.setRawText("Grace Hopper", ChangeOrigin.USER);

assertThat(model.problems().hasErrors()).isFalse();
assertThat(model.toResult().name()).isEqualTo("Grace Hopper");
\end{lstlisting}

Field-kind annotations choose generated model and view shapes. A
\api{@TextField} for \api{String}, \api{Integer}, \api{int},
\api{BigDecimal}, or \api{LocalDate} produces a \api{TextFieldModel} and a
\api{JTextField} view method. A \api{@DateField} currently uses a text-field
model with the built-in \api{LocalDate} converter but expresses date-specific
intent. A \api{@ComboBox} produces a \api{FieldModel} and a
\api{JComboBox<T>} view method; enum-valued combo boxes may also produce
generated option lists. A \api{@CheckBox} produces a \api{FieldModel} and a
\api{JCheckBox} view method. The generated view contract is narrow because it
exists only to expose the components the behavior plan needs.

The generated descriptor companion is the bridge from fields to surrounding
metadata. A descriptor can point to the field key, label resource, tooltip
resource, help topic, constraints, and display order. Dialog construction,
problem summaries, help routing, and tests can all ask the descriptor for the
same facts. A hand-written application can build equivalent descriptors, but
it is easy to forget one resource key or give one problem code a different
prefix. Generation removes that clerical risk for supported field kinds.

The generated view contract should remain small. If the behavior plan needs a
name field, the contract should ask for a name field. It should not expose the
whole panel, the layout manager, the parent dialog, or every helper label. A
small contract lets tests provide a minimal implementation and lets the
application replace the visual panel without rewriting generated bindings. It
also prevents generated code from reaching into view details that belong to
design and application code.

Generated bindings should be reversible. Installing text-field bindings,
validation tooltips, help behavior, and focus behavior creates listener
relationships. Those relationships need a lifecycle scope. A generated
\api{CustomerEditBehaviorPlan.install} method should add bindings to a scope
owned by the view or dialog. Closing that scope should detach document
listeners, value subscriptions, and input-map entries. Generation is useful
only if it removes listener sprawl without creating hidden listener leaks.

\begin{figure}[htbp]
\centering
\begin{minipage}[t]{0.48\linewidth}
\includegraphics[width=\linewidth]{\bindingFormFigure}
\end{minipage}\hfill
\begin{minipage}[t]{0.48\linewidth}
\includegraphics[width=\linewidth]{\dialogValidationFigure}
\end{minipage}
\caption{Generated form metadata feeds ordinary field models and Swing validation behavior. The generated view contract identifies components; the handwritten view still owns layout and density.}
\label{fig:generation-form-output}
\end{figure}

Figure~\ref{fig:generation-form-output} is deliberately visual. Generation
does not remove the need to design a screen. It gives the screen typed field
models, descriptors, validation metadata, and binding entry points. A
handwritten panel still arranges labels and editors with MigLayout. It still
chooses grouping, scrolling, command placement, and visual density. The
generated view contract lets the binding plan talk to the panel without
knowing the panel's concrete class.

Validation annotations are intentionally limited to common field-local rules.
\api{@Required} creates a required constraint and problem code. \api{@Length}
applies to string text fields. \api{@Regex} applies to string text fields and
requires a nonblank pattern. \api{@DecimalRange} applies to
\api{BigDecimal} text fields and requires at least one valid bound.
\api{@IntRange} applies to integer text fields. These constraints are good
generation inputs because they are local, declarative, and stable. A rule that
requires a server call, depends on a permission service, or compares several
records belongs in Java validation code.

Generated problem codes matter because validation messages often need stable
identity. A required-name problem may have a message resource, a help link, a
test assertion, and analytics. If the problem is only an English message, all
of those uses become fragile. If the problem has a generated
\api{ProblemCode}, the message can be localized and the test can still assert
the problem identity. The generated form model supplies the routine codes;
application validation can add richer problem codes when needed.

\api{@Help} on a generated form field contributes optional tooltip and help
topic metadata. The tooltip id is a resource identity. The help topic is a
stable help target. The generated descriptor can carry both, and Swing
behaviors can compose them with validation problems and disabled reasons.
Again, the generator does not open a help viewer. It supplies the durable
metadata that a help behavior and help service can use.

Generated options for enum combo boxes should be treated as convenient default
metadata, not as a complete option service. Enum declaration order is a fine
initial option order for small stable classifications. A dynamic option list
from a database, user permissions, or remote service remains application-owned.
The generated view contract can still expose a combo box, and the application
can still bind an observable option list where appropriate. Generation should
not turn dynamic product data into compile-time constants merely because
combo boxes exist.

The important habit is to read the generated model once. In the companion
example, the generated \api{CustomerEditFormModel} constructs text field
models from the original record, registers them with the form model, builds a
rule set for required and length validation, exposes descriptors, and
constructs a new \api{CustomerEdit} from parsed field values. There is no
reflection puzzle. The generated model is Java code that a careful developer
could have written by hand and would not want to write repeatedly.

\section{Generated tables from records}

\api{@SwingTable} marks a record as a generated table row source. Record
components annotated with \api{@Column} become table columns. Unannotated
components remain part of the row record but are ignored by table metadata
generation. This is useful when a row carries an internal id, cached status,
or detail object that should not appear as a visible column. Generation
follows the annotations, not every property reachable from the row.

\begin{lstlisting}[caption={Annotated table source used by the companion example.}]
@SwingTable(id = "book/customer/table")
public record CustomerRow(
        @Column(header = "customer/name", width = 180)
        String name,

        @Column(header = "customer/status", width = 90)
        String status) {
}
\end{lstlisting}

The generated \api{CustomerRowColumns} companion contains a table key, column
keys, \api{Column} instances, and \api{ColumnDescriptor} instances. The
generated \api{CustomerRowTableResources} companion contains header and
tooltip resource keys. The generated \api{CustomerRowTableDescriptor}
contains a \api{TableDescriptor<CustomerRow>} and a factory for an
\api{ObservableTableModel}. The generated bindings companion installs table
models and selection bindings into Swing views. The table chapter explains
these runtime classes in detail; here the point is how generation gives them
stable source.

This generated table surface solves the old \api{JTable} column-index problem.
In handwritten Swing, column zero may mean name in one method, status in
another method, and id in a renderer after a later edit. A generated
\api{ColumnKey<CustomerRow,String>} names the column independently from its
current view index. A generated descriptor can preserve record-component order
as the default order while table state, sorting, filtering, and visibility
operate with stable column identity. The code still has to respect view and
model coordinates, but it no longer needs to guess what an integer column
means.

Generated column value extraction is also a performance boundary. A descriptor
can call a record accessor directly. The table model does not need reflective
property lookup for every cell. That matters in large tables, especially when
sorting, filtering, repainting, and tooltips all ask for cell values. Direct
generated accessors are simple Java calls that JIT compilation can understand.
The generator supplies the repetitive adapter code once instead of making each
table author handwrite a switch statement.

Column metadata is richer than a header. The \api{@Column} annotation can
declare a column id, header resource id, optional tooltip resource id,
persistence id, width, minimum width, maximum width, order, initial
visibility, sort capability, filter capability, hideability, and editability.
Defaults are derived from the table id and record component name. The defaults
are helpful, but they are not free to change once persisted state or resource
bundles depend on them.

Table state is where stable ids become concrete. A user may resize, reorder,
hide, show, and sort columns. The table-state controller can persist those
preferences by table id and column persistence id. If a display header changes
from \api{"Customer"} to \api{"Name"}, saved state should still restore the
same column. If the generated column id changes accidentally, saved state may
be lost or applied to the wrong column. This is why column ids and
persistence ids are generated metadata, not visible labels.

\begin{migrationbox}{Rename the Resource, Not the Column}
If a table header changes for copyediting or localization, change the header
resource. Change a generated column id only when the column's durable identity
has changed and a table-state migration has been considered.
\end{migrationbox}

Generated table descriptors also help with large-data discipline. The
descriptor says which column extracts which value and which columns are
sortable, filterable, editable, or hideable. An \api{ObservableTableModel}
created from that descriptor can report precise cell and row changes when the
observable row source changes. A generated descriptor does not make a million
rows cheap by itself, but it prevents a common source of table slowness:
index-based column code scattered through renderers, sorters, and export
helpers.

Editable generated columns require special caution. A record is immutable, so
an editable column must create an updated row value, usually by calling the
record constructor with one component changed. The processor can generate
helpers for supported editable columns, but the application still owns the
policy around editing. Is the edit committed immediately? Does it mark the
row dirty? Does it require validation? Does it update a remote source? Does it
preserve selection? Generated code can provide the mechanical updater; the
workflow remains application code.

Selection binding is another reason to generate table metadata. A table
selection can be represented as selected view row, selected model row, selected
row object, selected row id, or selected column. Handwritten screens often mix
these. A generated table binding can install a model and bind selection
through the standard table-selection helpers. The presenter can then observe a
\api{WritableValue<CustomerRow>} or row id value instead of reading
\api{JTable.getSelectedRow()} in every command. Generated metadata does not
remove Swing's coordinate systems, but it gives the application a consistent
place to cross them.

Filtering and sorting should respect generated column capabilities. A column
marked not sortable should not receive a sorter toggle merely because it is
visible. A column marked not filterable should not appear in a filter builder
unless application policy overrides the descriptor. Generated descriptors make
these capabilities reviewable. They also make table preference dialogs and
column visibility menus easier to build because the metadata is typed and
ordered.

Generated table resources should be audited before table screenshots are
taken. Header resources affect width, wrapping, tooltip text, accessibility
names, export labels, and screenshot diffs. If a generated header key is
missing, a table may still render something, but the screenshot will review a
fallback rather than the intended UI. A table-resource test can iterate the
generated resources companion and require all headers. This is cheap compared
with debugging a screenshot that silently used placeholder text.

A table descriptor is also a documentation object. A developer reading the
generated descriptor can see which columns exist, their order, widths, value
types, editability, and resource keys. In a rich client with dozens of
business tables, that generated descriptor is often easier to inspect than the
final Swing table after column state, renderers, and filters have been
installed. The descriptor is the table's declared contract; the \api{JTable}
is one presentation of that contract.

Tooltips and help topics on columns deserve the same separation as form fields.
A column header tooltip is a resource. A help topic is a durable help target.
A cell problem is a problem with a row and column target. A renderer may show
a visual marker. A table-cell tooltip binding may compose the problem, static
help, and F1 hint. Generated column metadata gives each concern a stable
column identity. It does not collapse them into one text string.

Generated table bindings should be read as installation helpers, not as a
replacement for table design. A table can still use custom renderers, row
sorters, Glazed Lists integration, selection models, column visibility menus,
and state persistence. The generated binding can install the model and
selection wiring. The application then adds specialized visual and behavioral
policy. This is the same division as generated forms: generated code owns
repeated contracts, not the whole screen.

The companion example's generated \api{CustomerRowTableDescriptor} is a good
file to open. It constructs a \api{TableDescriptor} from generated columns and
offers a \api{tableModel} factory for an observable row source. This is boring
code. That is the point. If every table in the application needs the same
shape of descriptor and model factory, the processor should write it once
consistently rather than asking developers to type and review the same pattern
again and again.

\section{Generated actions and command metadata}

\api{@UiAction} marks a no-argument \api{void} method as a generated UI action
source. The processor validates the method shape and writes an owner-specific
actions companion. For each action it generates an \api{ActionKey}, text
resource key, optional tooltip resource key, \api{ActionDescriptor}, command
factory, and command-set factory. The annotated method remains application
code. The generated companion supplies the stable command metadata and the
mechanical factory path.

\begin{lstlisting}[caption={Annotated presenter actions.}]
final class CustomerPresenter {
    @UiAction(
            id = "customer/save",
            text = "customer/save/text",
            tooltip = "customer/save/tooltip",
            mnemonic = KeyEvent.VK_S,
            acceleratorKey = KeyEvent.VK_S,
            acceleratorModifiers = InputEvent.CTRL_DOWN_MASK)
    void save() {
        service.save(form.toResult());
    }

    @UiAction(id = "customer/active", selectable = true)
    void toggleActive() {
        activePreview.setValue(!activePreview.value(), ChangeOrigin.USER);
    }
}
\end{lstlisting}

The generated descriptor is an \api{ActionDescriptor}. The generated command
factory returns a \api{MutableCommand}. The generated command-set factory
returns a \api{CommandSet}. Swing views can adapt those commands through
\api{SwingActionAdapter} or generated command behaviors. Tests can require the
command by generated action key. Resource tests can require generated command
text keys. Menu construction can iterate generated descriptor lists. A
shortcut can be installed from the generated accelerator metadata. One
annotated method thus becomes a coherent operation identity across many Swing
surfaces.

The generated action companion should be read as operation metadata, not as a
presenter replacement. The presenter still owns state such as dirty flags,
validation problems, selected rows, permissions, and busy values. The
generated command factory can create a mutable command that calls the annotated
method. Application code can then wire enablement and selected state from the
presenter's observable values. This preserves the command chapter's rule: the
operation has stable identity, but runtime availability is live state.

Generated factories are useful even when a screen needs custom command
construction. A presenter may ignore the generated factory and use the
generated descriptor with a custom handler that starts a task, wraps errors,
or records telemetry. That is still a good use of generation. The descriptor,
keys, resources, and action id remain compiler-visible and shared across
surfaces. Generation should not force every command through one factory style
when application workflow needs a richer command object.

Command-set factories help with review. A generated \api{commands(owner)}
method can expose every annotated action in deterministic order. A menu
builder, toolbar builder, command palette, or test can inspect the same set.
Hidden listeners become easier to spot because they do not appear in the
generated action vocabulary. A reviewer can ask why a visible button is not
backed by a generated or handwritten command descriptor. This question often
finds old Swing code that bypasses enablement and shortcut policy.

Selectable actions deserve a clear ownership model. The \api{selectable} flag
marks toggle-style action metadata. The generated factory should use a toggle
command path, and Swing toggle controls should observe the command's selected
value. The annotation does not decide where the selected state ultimately
comes from. A presenter may keep it in a model value, derive it from current
view state, or update it in the handler. What generation supplies is the
metadata that this command has selected state at all. That prevents a toolbar
toggle and check-box menu item from being generated as ordinary one-shot
actions by accident.

Descriptors are not buttons. Generated action text and tooltip ids are
resources, not visible strings. Mnemonic and accelerator values are metadata,
not installed key bindings. The Swing integration layer decides where a
descriptor is presented: menu, toolbar, popup menu, button row, command
palette, or root-pane key binding. This division is important because the same
operation may have several surfaces with different density and focus policy.
Generation names the operation; the view presents it.

The processor currently rejects methods with parameters and methods that do
not return \api{void}. That may seem strict, but it keeps generated command
factories direct. A generated factory can call \api{owner.save()} without
inventing an argument source. If an operation requires a selected row, current
form state, or application service, the presenter method can read those facts
from presenter state. If an operation needs a parameter supplied by the caller,
it is probably not a simple UI action command and should be modeled more
explicitly.

Duplicate action ids are rejected within one owner. This protects command
sets, menu descriptors, diagnostics, tests, and resource keys. If two methods
have the same action id, either one of them is a duplicate presentation of the
same operation and should not have its own generated action, or they are
different operations and need different ids. The processor should not guess.

Accelerator metadata has its own validation. A nonzero accelerator modifier
mask without an accelerator key is invalid. This catches source declarations
that cannot produce a meaningful key stroke. Platform-specific accelerator
policy may still override generated defaults during installation. The
descriptor supplies a stable default; the view or application policy decides
where and whether to install it.

Generated action companions commonly expose descriptor lists for menus and
toolbars. Treat those lists as defaults, not immutable product design. A menu
may include all descriptors. A toolbar may show only frequent operations. A
popup menu may include context-sensitive actions after selection refresh. A
command palette may include operations that have no visible button. The same
generated descriptors can feed several policies because they are operation
metadata rather than component instances.

The \generatedActionExample{} companion demonstrates generated action metadata
on a small source class. It is worth reading together with the command chapter.
The generation chapter explains how the descriptor appears. The command
chapter explains how the descriptor becomes runtime command state and Swing
actions. The important habit is not to stop at the generated descriptor. A
real screen still needs command enablement, selected-state ownership, task
handoff, error policy, and lifecycle cleanup.

The same action id may appear in resources, tests, documentation, telemetry,
and persisted shortcut preferences. This is why \api{@UiAction(id = "...")}
must be treated as more than a method decoration. Renaming \api{save()} to
\api{saveCustomer()} is a Java refactoring. Renaming \api{customer/save} is a
contract change. If a generated action id is already used outside the source
file, migrate it deliberately or keep the id stable while refactoring the
method name.

Generated action descriptors are also useful for accessibility. If a toolbar
uses icon-only buttons, the generated text resource can supply an accessible
name and the tooltip resource can help supply an accessible description. If a
menu uses localized text, the action id still identifies the operation. If a
shortcut is shown in help, the descriptor provides the default accelerator.
Accessibility metadata should not be recreated independently for each
presentation of the same command.

\section{Build integration and generated source}

Annotation processing is a build concern. The annotations belong on the
compile classpath because source code imports them. The processor belongs on
the annotation-processor path because javac invokes it during compilation.
Runtime code should not need the processor jar to start the application. This
separation keeps generated metadata available at compile time without turning
the processor into runtime infrastructure.

\begin{lstlisting}[caption={Typical Gradle dependency shape.}]
dependencies {
    implementation project(":corusco-core")
    implementation project(":corusco-swing")
    implementation project(":corusco-annotations")
    annotationProcessor project(":corusco-processor")
}
\end{lstlisting}

The exact module names and dependency scopes may differ in a project, but the
shape should remain clear. The application compiles against core, Swing, and
annotation APIs. The compiler invokes the processor. The generated sources are
compiled into the same source set as the annotated source. At runtime, the
application uses generated classes and Corusco runtime libraries, not the
processor.

Generated source directories should be visible to the IDE and marked as
generated. Developers should be able to navigate to a generated class from an
import, set a breakpoint if debugging generated behavior is useful, and read
the source in code review. They should not edit generated files. If a
generated file is wrong, edit the annotated source or the processor. Manual
changes to generated output are build artifacts that will disappear or, worse,
mislead the next reader.

Stale generated source is a build smell. If an annotation is removed and the
old generated class remains visible, clean the build and inspect source-set
configuration. Gradle's Java compilation tasks should treat generated sources
as task output so obsolete files do not masquerade as current contracts. Stale
output can be especially confusing with generated companions because their
names are predictable. A class named \api{CustomerEditFields} may still compile
even after the annotated form source changed. A clean build is the first test.

In modular applications, generated classes belong to the module that owns the
annotated source. They do not automatically require public exports. A view
package may use generated companions internally. A test module may require
access to those packages. A public API module should not export generated UI
helpers merely because they exist. The module boundary should reflect the
application contract, not the presence of generated files.

Annotation processing and jlink packaging should remain boring. The runtime
image needs the compiled generated classes and runtime Corusco modules. It
does not need javac or the annotation processor. If a packaged application
tries to load processor classes at startup, something has crossed the
compile-time boundary incorrectly. The packaging chapter returns to runtime
images; the generation-specific point is simple: generated source is compiled
away into ordinary classes before packaging.

Resource files are adjacent but separate. Generated resource keys tell the
application which resources are expected. They do not write localized text
into a bundle unless a separate tooling step chooses to do that. A build may
include a resource audit that iterates generated resource keys and verifies
that required bundles contain values. That audit should run before screenshot
generation, because missing labels make visual review noisy and hard to
interpret.

Multi-module builds need an explicit ownership rule. If a form source lives in
one module and a Swing view lives in another, the generated companions must be
on the consuming module's compile path. That may mean moving the annotated
source to a shared UI-model module, moving the view into the same module, or
exporting a generated package intentionally. Do not solve this by copying
generated files between modules. Generated files are outputs of their source
module. Copying them creates a stale-shadow problem that the compiler cannot
explain well.

IDE configuration should make generated code visible but visually distinct.
Developers need completion for generated companions and navigation from
generated constants to usages. They also need to know that editing a generated
file is the wrong fix. Most IDEs mark generated source roots differently. Use
that feature. A generated file that looks like handwritten source invites
manual edits; a generated file that cannot be navigated invites distrust.

Continuous integration should include at least one clean build path. Incremental
builds are valuable during development, but generation bugs and stale-output
bugs often hide in incremental state. A clean CI build proves that annotated
sources alone are sufficient to recreate the generated companions. If clean CI
and incremental local builds disagree, the generated-source task boundaries
need attention before application code is blamed.

Publishing generated artifacts requires a deliberate policy. An application
module normally publishes compiled classes that include generated companions.
A library module may expose generated companions as part of its API if other
modules are expected to use them. That decision makes generated names and ids
more compatibility-sensitive. If generated companions are internal, keep them
in non-exported packages. If they are public, document them like public API.

Generated source should not be committed by default unless the project has a
specific reason. Committing generated files can make diffs noisy and create
merge conflicts. Not committing them requires reliable processor configuration
and generated-source visibility. Some projects choose to publish generated
source jars for debugging; that is different from treating generated output as
hand-maintained source. The rule is not universal, but the policy should be
explicit.

Generated sources are also useful for documentation and examples. Companion
examples should compile annotated sources through the real processor. The text
can then refer to \generationExample{} as an anchor without copying a long
generated class into prose. If a listing needs to show generated shape, it
should show a small excerpt and tell the reader which example produces the
full source. This prevents documentation from becoming a second, stale copy of
generated output.

\begin{examplebox}{Generation}
\generationExample{} is the navigation anchor for the generated customer
example. The annotated records \api{CustomerEdit} and
\api{CustomerRow} live in the same package. Their generated companions are
emitted next to those source types during the \api{corusco-examples} build.
\end{examplebox}

When upgrading Corusco, generated source is a review artifact. A processor
change may alter generated names, ids, descriptors, import shape, or behavior
installation. Some changes are improvements. Some are compatibility breaks.
Review representative generated output before accepting a processor upgrade in
a large application. A generated-source test can lock down the project-specific
contracts that matter most.

\section{Diagnostics and generated-source tests}

Good generation fails early and locally. If \api{@SwingForm} is applied to a
class instead of a record, the processor should report that the annotation is
supported only on records. If a record component has both \api{@TextField} and
\api{@CheckBox}, the processor should report that only one field-kind
annotation may be present. If \api{@Length} is placed on a
\api{BigDecimal} text field, the processor should report that length applies
to string text fields. The error should point at the source declaration the
developer can fix, not at a generated file that failed later.

Diagnostics are part of the generated API because they teach the rules. A
developer encountering generation for the first time will learn from processor
messages as much as from documentation. Messages should use the annotation
names, source type names, conflicting ids, and expected constraints. A message
that says \api{"invalid metadata"} wastes the compiler's opportunity. A
message that says \api{"Duplicate @Column id in CustomerEdit:
customer/search/name"} tells the developer exactly which contract is broken.

The processor test suite should contain both success and failure cases.
Success cases prove that representative annotated sources generate the
expected companions. Failure cases prove that invalid source is rejected with
useful messages. These are not merely processor tests. They are examples of
the contract that application teams rely on. When the processor changes, the
tests say which generated shapes and diagnostics are intentionally stable.

\begin{lstlisting}[caption={Generated-source compiler tests exercise the real javac path.}]
GeneratedSourceCompilation result =
        GeneratedSourceCompiler.in(tempDir)
                .withProcessor(new CoruscoAnnotationProcessor())
                .compile("demo/CustomerEdit.java", source);

assertThat(result.success()).isTrue();
result.assertGeneratedSourceContains(
        "demo/CustomerEditFields.java",
        "TextFieldKey.of(\"customer/name\"");
\end{lstlisting}

\api{GeneratedSourceCompiler} exists for this kind of focused test. It writes
source text into a temporary compilation directory, runs javac with the
provided processors, stores generated sources under a known generated-output
directory, and returns success status plus readable diagnostics. The companion
\api{GeneratedSourceCompilation} can return generated source text or assert
that generated source contains important snippets. This is more realistic than
calling processor internals directly because it exercises javac, the
annotation-processing API, source paths, generated output, and diagnostics
together.

Generated-source tests should be small. A test for duplicate table column ids
does not need a complete application. It needs one annotated record with two
columns declaring the same id. A test for action parameter rejection needs one
annotated method with a parameter. A test for generated form view contracts
needs one record with representative field kinds. Small samples make
diagnostics clear and keep processor tests fast enough to run frequently.

There is a trap in testing generated output too literally. The exact import
order, whitespace, or line wrapping of generated source is rarely the public
contract. Important constants, type names, descriptor calls, method names,
diagnostic messages, and key ids are often the contract. Generated-source
tests should assert the parts that matter. If every formatting detail is
locked, refactoring the generator becomes unnecessarily expensive.

Application tests should complement processor tests. A processor test can
prove that \api{CustomerEditFormModel} contains a text field model for
\api{name}. A core form test can instantiate the generated model and verify
parse and validation behavior. A Swing test can implement the generated view
contract and install the generated behavior plan. A resource test can require
generated resource keys. A screenshot test can use the same generated view
contract and resources to render a stable screen. Each layer tests a different
boundary.

Generated behavior plans deserve lifecycle tests. Create a minimal view
implementation, a generated form model, and a \api{BehaviorScope}. Install the
generated plan, change the document text, and assert that the model changes.
Close the scope, change the document again, and assert that the disposed
binding no longer updates the model. This test is small, but it protects the
promise that generation removes listener sprawl without leaving listeners
behind. It is also much faster than testing the same fact through a full
dialog workflow.

Generated action companions deserve command tests. Create an owner object with
annotated methods, build the generated command set, require a command by
generated action key, and execute it. Then test descriptor metadata separately:
text resource id, tooltip id, selectable flag, mnemonic, accelerator, and
descriptor ordering. The execution test proves the factory path. The metadata
test proves the generated command contract. Mixing all of this into a button
click test makes failures harder to diagnose.

Generated table descriptors deserve model tests. Build an observable row list,
create the generated table model, and assert column count, names, value
extraction, column classes, editability, and precise update events. If a column
is editable, test the generated row updater or the application policy around
it. If table state persistence uses generated persistence ids, test that those
ids are the expected stable values. A screenshot can then focus on visual
rendering rather than discovering that the descriptor is wrong.

Generated resource tests should be data-driven. A generated resources
companion often contains constants that can be listed directly or through a
small project convention. The test should require all mandatory labels and
headers in the sample resource set. Optional tooltips may be resolved with
fallbacks, but a missing required command text should fail before the screen is
opened. This is especially important for visual examples, because screenshots
should show deliberate text rather than missing-resource markers.

\begin{figure}[htbp]
\centering
\begin{minipage}[t]{0.48\linewidth}
\includegraphics[width=\linewidth]{\testingKeysFigure}
\end{minipage}\hfill
\begin{minipage}[t]{0.48\linewidth}
\includegraphics[width=\linewidth]{\testingHarnessFigure}
\end{minipage}
\caption{Generated keys are useful in tests because they separate the mechanics of finding components from the visual review of rendered text and layout.}
\label{fig:generation-testing-output}
\end{figure}

Figure~\ref{fig:generation-testing-output} connects generated metadata to the
testing chapter. A generated view contract can give a test typed component
handles. A screenshot harness can render the real Swing panel with FlatLaf and
MigLayout. The harness should find components by generated keys or view
contract methods, while the screenshot itself shows real localized text. This
keeps test mechanics stable when text changes and keeps visual review honest
when text does change.

Diagnostics should also protect stable-id grammar. Invalid ids should fail at
compile time, not later when a resource lookup fails or table state cannot be
restored. The processor should reject blank ids, duplicate ids in a local
scope, invalid tooltip conflicts, invalid width ranges, invalid validation
bounds, unsupported field kinds, and inconsistent accelerator metadata. A
strict processor may feel less forgiving during the first compile, but it is
more humane than letting a malformed screen reach runtime.

\section{Review, migration, and versioning}

Generated code changes the shape of review. A reviewer should inspect three
things: the annotated source, the generated source contract, and the
handwritten code that consumes the generated companions. The annotated source
answers whether the stable facts are correct. The generated source answers
whether the processor interpreted them as expected. The handwritten code
answers whether the screen still owns layout, workflow, and application
policy. Skipping any one of these views makes generation either too trusted or
too feared.

The generated naming map should be memorized by the team. It is small enough:
\api{CustomerEditFields} for field keys, \api{CustomerEditResources} for form
resources, \api{CustomerEditProblems} for problem codes,
\api{CustomerEditDescriptors} for field descriptors,
\api{CustomerEditFormModel} for editable model state, \api{CustomerEditView}
for the Swing component contract, \api{CustomerEditBehaviorPlan} for direct
behavior installation, and \api{CustomerEditBindings} for the binding facade.
Tables and actions have similarly predictable names. Predictability is not a
cosmetic feature. It is what lets a developer navigate the generated surface
without consulting a manual on every change.

\begin{tabularx}{0.96\linewidth}{@{}p{0.30\linewidth}X@{}}
\toprule
Generated area & Review questions\\
\midrule
Fields and resources & Are ids stable, readable, and independent from visible text?\\
Problems and validation & Are generated constraints field-local, and are richer rules still in Java?\\
View contracts & Does the handwritten view own layout while satisfying the generated contract?\\
Behavior plans & Are installed behaviors reversible and scoped to the view lifecycle?\\
Table descriptors & Are column ids and persistence ids durable enough for saved state?\\
Action descriptors & Do buttons, menus, toolbars, and shortcuts share command identity?\\
Build setup & Are annotations on the compile path and processors on the processor path?\\
\bottomrule
\end{tabularx}

Migration should be incremental. Start with a form whose handwritten metadata
is already repetitive. Annotate the record, generate the form model and keys,
and keep the existing panel layout. Add tests for generated model behavior and
resource presence. Then install generated behavior into the panel. After that,
migrate table descriptors or action descriptors. Each slice should reduce one
class of duplication while preserving the user's visible workflow. A
big-bang migration makes it hard to tell whether a failure is a processor bug,
a changed id, a changed resource, a changed layout, or a changed workflow.

Do not migrate only the annotations. A record with \api{@SwingForm} is useful
when application code uses the generated companions. If the panel still
creates unrelated field keys by hand, tests still search for labels, and
validation still emits raw strings, generation has become decoration. The
payoff comes when fields, resources, problems, behavior, and tests share the
generated identity.

Do not migrate layout into generation by accident. When a team first sees a
generated view contract, it is tempting to ask the generator to produce the
panel too. Resist that as the default. Layout is where product design, user
workflow, density, grouping, and accessibility meet. MigLayout makes ordinary
handwritten layout concise enough to review. Generated code can require a
\api{nameField()} method and install a binding. It should not decide that the
name field belongs in row three, column two of every application dialog.

Versioning belongs in the migration plan. If a generated column id changes,
write or consider a table-state migration. If an action id changes, review
shortcuts, command palettes, help topics, and tests. If a form field id
changes, review resources, problem messages, screenshot fixtures, and
component keys. If a generated class name changes because the source type was
renamed, decide whether external tests or examples refer to it. The compiler
will catch some references. It will not migrate saved table state or external
documentation by itself.

Generated source can make deprecation easier. A processor may keep an old
resource key constant for a release while introducing a new one, or a table
state migration may map old persistence ids to new ids. Such compatibility
features should be explicit and temporary. Silent aliasing forever makes the
generated surface harder to reason about. The usual rule applies: stable ids
are durable; changed ids are migrations; migrations should be named.

Resource authors should be included in the contract. Generated resource keys
are not merely Java constants. They are work items for localization and UX
copy. A resource audit can list generated keys missing from a locale. A
pseudo-localization screenshot run can show whether generated labels and table
headers still fit. A generated key naming scheme that resource authors can
read will reduce mistakes. A scheme full of Java implementation names will
spread Swing internals into localization work.

Generated code should be read more often during processor development than
during everyday application work. In normal screen work, a developer may read
the generated companions for a new form or table and then trust the pattern.
When changing the processor, read many generated companions. Check small
forms, enum combo boxes, validation combinations, editable tables, action
metadata, and module boundaries. The processor is a multiplication tool; a
small bad assumption can appear in every generated screen.

Review generated diffs by intent. A processor refactoring that only changes
import ordering should not be described as a UI behavior change. A processor
change that renames generated constants is a source compatibility change. A
change that alters generated ids is a resource, test, and persistence
compatibility change. A change that alters behavior-plan installation order
may be a runtime behavior change. Sorting generated diffs into these
categories keeps review focused on consequences rather than line count.

Generated naming should avoid clever compression. A file named
\api{CustomerEditBinds} is shorter than \api{CustomerEditBindings}, but the
longer name is clearer and matches the vocabulary used elsewhere. Generated
code is read in moments of uncertainty: when a build fails, a binding behaves
strangely, or a test cannot find a key. Clear generated names reduce the cost
of those moments. They also make examples and documentation easier to write
without long explanations.

Migration from handwritten metadata should leave bridges only temporarily. It
is acceptable to introduce generated field keys while an old presenter still
uses handwritten problem codes. It is risky to keep both systems forever. Two
parallel identity systems mean two places for bugs. Plan the migration slice,
add tests, move consumers to generated companions, and remove obsolete
handwritten constants when the slice is complete. Generated code should
replace repetition, not add a second layer beside it.

When a processor limitation blocks a migration, prefer an explicit handwritten
bridge to an annotation contortion. If a field kind is not supported, declare
the field key and descriptor by hand for that field. If a table column needs a
special renderer, keep the generated descriptor for identity and install the
renderer explicitly. If an action needs arguments, write a custom command
using generated resource keys where useful. Such bridges are honest. Forcing a
limited annotation to mean something it does not support creates future
confusion.

Documentation should name both the annotation and the generated result. Saying
``add \api{@SwingForm}'' is not enough. Say that the source will generate a
form model, field keys, descriptors, resources, a view contract, a behavior
plan, and bindings facade. Saying ``add \api{@UiAction}'' is not enough. Say
that it generates an action key, descriptor, resource keys, command factories,
and descriptor lists. Readers learn generation best when the input and output
are always paired.

\section{Boundaries that should remain handwritten}

It is worth naming what generation should not own, because otherwise every
annoying repetition begins to look like an annotation request. The first
boundary is layout. A generated view contract can require methods for
components. A generated behavior plan can bind them. The panel layout should
still be reviewed as Swing code. A business screen with dense forms, nested
tables, split panes, command rows, and validation summaries needs design
judgment. An annotation on a record component cannot see the whole screen.

The second boundary is service workflow. A generated action factory can call
\api{owner.save()}. It should not know which repository is used, whether a
transaction spans several aggregates, whether a confirmation dialog appears,
or whether a virtual thread task is started. Those decisions need application
context and tests. Generation gives the action stable identity and a command
factory; the presenter owns the workflow.

The third boundary is cross-field and cross-record validation. Field-local
constraints such as required text and length are good generation inputs.
Validation that compares start and end dates, calls a server, checks
permissions, or depends on a selected row is ordinary Java validation. It can
still emit generated problem codes or target generated field keys. It should
not be squeezed into annotation attributes merely to make a form look more
generated.

The fourth boundary is dynamic choice data. A generated enum option list is
acceptable because enum constants are compile-time facts. A list of customers,
roles, countries, project templates, or remote statuses is not the same kind
of fact. Such data may need loading, filtering, permissions, localization, and
refresh. Generate the field identity and view contract; bind the dynamic
options through application state.

The fifth boundary is rendering policy. A generated table column can have a
value type, width, resource keys, and editability. It should not decide the
custom renderer for a GIS geometry, a financial sparkline, or a color-coded
risk cell unless the project deliberately adds a narrow generated renderer
contract. Renderers often require visual and performance review. Generated
metadata can supply column identity and value extraction while renderers stay
explicit.

The sixth boundary is security. A generated command descriptor can name the
Delete operation. It should not decide whether the current user may delete the
selected customer. Permission checks belong in presenters, command enablement,
services, or policy objects. Generated ids can help diagnostics and help
messages explain permission decisions, but the authority should be explicit
application code.

The seventh boundary is long-running work. A generated command factory may
call a presenter method that starts a task. It should not choose the task
service, cancellation policy, stale-result strategy, progress model, or busy
overlay. Those topics belong to the lifecycle and background-work chapters.
Generation can reduce wiring around commands and fields. It should not make
concurrency implicit.

These boundaries are not a limitation of the processor. They are what keep
generated code useful. A generator that tries to own layout, services,
security, rendering, dynamic data, and concurrency becomes a second
application framework hidden behind annotations. Corusco's better role is
smaller and sharper: generate stable identity, descriptors, models, and
routine behavior entry points; let application code express the screen.

The boundary can be tested in code review with a simple question: could this
decision be made correctly by looking only at the annotated source element?
If the answer is yes, generation may be appropriate. A record component's
value type and field id can be known from the component. A table column's
default width and resource id can be declared on the component. An action's
default shortcut can be declared on the method. If the answer is no because
the decision depends on several fields, current user, selected row, loaded
data, application mode, or visual grouping, the decision should remain in
ordinary code.

This question prevents both underuse and overuse. Underuse appears when every
field key, resource key, problem code, and descriptor is handwritten even
though the record declaration already says the necessary facts. Overuse
appears when annotations begin to encode workflow, layout, permissions, or
service calls. The useful middle is broad enough: stable keys, descriptors,
routine models, routine table metadata, and routine command metadata. That is
where Corusco generation should be strong.

\section{Worked review of the companion example}

Open \generationExample{} as a navigation anchor, then open the two annotated
records in the same package. \api{CustomerEdit} is a generated form source.
\api{CustomerRow} is a generated table source. The source code is intentionally
small, because the generated companions are the thing to inspect. After the
examples module compiles, the generated-source directory contains companions
beside those records. The generated files are not meant to be edited. They are
meant to be read.

The form source starts with \api{@SwingForm(id = "book/customer/edit")}. That
id becomes the prefix for generated field ids and problem codes. The
\api{name} component is a text field with required and length constraints.
The \api{status} component is a text field without extra validation. From
those facts, the processor emits two text field keys, descriptor constants,
resource keys, problem codes for the name constraints, a form model, a view
contract, behavior plan, and bindings facade. None of that requires runtime
reflection over the record.

The generated \api{CustomerEditFields} file is the simplest inspection point.
It contains \api{NAME} and \api{STATUS} constants. Their ids are
\api{book/customer/edit/name} and \api{book/customer/edit/status}. Their owner
type is \api{CustomerEdit}. Their value type is \api{String}. A reader can
now predict the generated descriptors and problem targets. The file is short
and boring. This is exactly what generated key companions should be.

The generated \api{CustomerEditFormModel} is more instructive. Its constructor
requires an original \api{CustomerEdit}. It creates text field models using
the generated keys and string converters. It registers those models with the
abstract form infrastructure. It builds a rule set that validates required and
length constraints for the name field. It exposes descriptors in record
component order. Its \api{createResult} method constructs a new
\api{CustomerEdit} from parsed field values. Reading that class explains how
generated field metadata becomes a mutable edit model.

The generated \api{CustomerEditView} is a contract, not a panel. It requires
methods that return the Swing components used by the behavior plan. A
handwritten panel may use MigLayout and FlatLaf, add labels and command rows,
wrap fields in a scroll pane, or compose the form into a larger workspace. As
long as it implements the generated contract, the generated behavior plan can
install text-field bindings and validation behavior. The panel remains real
Swing code.

The table source starts with \api{@SwingTable(id = "book/customer/table")}.
The \api{name} column declares a header resource id and width. The
\api{status} column does the same. The generated columns companion contains
the table key, column keys, descriptors, and column instances. The generated
resources companion contains header keys. The generated table descriptor
constructs a \api{TableDescriptor} from the generated columns and provides an
\api{ObservableTableModel} factory. The generated bindings companion can
install the model and bind selection.

Now perform the compatibility review. If the visible header text should change,
change the resource value associated with \api{customer/name} or choose a
better header resource id. If the column should keep saved width and order,
do not change the column id or persistence id casually. If the status field
becomes a combo box in the form, update the field-kind annotation and inspect
the generated model, view contract, and options. If the screen needs a new
Refresh command, add \api{@UiAction} to the presenter and inspect the
generated action descriptor.

Finally, decide what the example does not show. It does not show generated
layout. It does not show service workflow. It does not show a real resource
bundle. It does not show table-state migration. It does not show a dynamic
combo-box option source. Those omissions are not accidental. The example is a
generation example: annotated source in, generated companions out. The other
chapters show how those companions participate in real screens.

\section{Exercises and checklist}

\begin{exercisebox}
\begin{enumerate}
\item Annotate a small record with \api{@SwingForm}. Build the module and list
every generated companion by name.
\item Open the generated form model and identify which lines construct field
models, which lines build validation rules, and which lines create the result
record.
\item Add \api{@Help(topic = "...")} to a generated form field. Find the help
topic in the generated descriptor and explain how F1 behavior would use it.
\item Annotate a table row record with two columns. Change only the header
resource id and explain why saved table state should not be affected.
\item Create an invalid table source with duplicate \api{@Column} ids and
write a generated-source compiler test that asserts the diagnostic.
\item Add \api{@UiAction} to a no-argument presenter method. Find the generated
action key, descriptor, command factory, and command-set factory.
\item Decide whether a cross-field validation rule belongs in an annotation or
in Java validation code. Justify the boundary.
\item Run a clean build and verify that obsolete generated companions disappear
after an annotation is removed.
\item Inspect a generated view contract and implement it in a MigLayout panel
without letting generated code own layout.
\item Review one generated id as if it were already referenced by resources,
tests, help, and persisted UI state. Decide whether renaming it is safe.
\end{enumerate}
\end{exercisebox}

\begin{center}
\textbf{Generation checklist.}
\end{center}

\begin{itemize}
\item Put annotations on stable source facts, not transient workflow choices.
\item Treat generated ids as compatibility-sensitive once screens are used.
\item Keep annotations on the compile classpath and processors on the
annotation-processor path.
\item Keep generated source visible in the IDE and marked as generated.
\item Read generated companions during review when annotations change.
\item Use generated keys and descriptors instead of reconstructing ids by hand.
\item Keep layout, service workflow, security, concurrency, dynamic options,
and specialized rendering in handwritten code unless a narrow generated
contract is deliberately designed.
\item Test processor contracts with generated-source compiler tests.
\item Test generated form models without Swing and generated behavior plans
with EDT Swing tests.
\item Audit generated resource keys before screenshot generation.
\item Migrate incrementally: form metadata, table metadata, and action
metadata can move one slice at a time.
\item Clean the build when generated source appears stale.
\end{itemize}

The proper attitude toward generated code is neither suspicion nor reverence.
It is source code produced by a tool the project controls. Read it when the
contract matters. Test it when the contract must remain stable. Improve the
processor when many generated files need the same change. Keep handwritten
code where judgment belongs. Done this way, compile-time generation gives
Swing applications a rare advantage: the repetitive contracts are checked by
the compiler, while the screen remains ordinary Java that a developer can
understand.
